\section{The Poisson probe and its Euler product}\label{sec:poisson}

We insert the theta sum \(B_m(Z)\) of
\Cref{eq:reflection-B-definition} into an additive average.  The reflected
second moment bounds this average from above.  Poisson summation will give a
second expression whose principal frequency contains the reciprocal of the
target \(L\)-function.  The use of an auxiliary additive polynomial while
retaining the unrestricted cube factors follows the method of
\cite[Part~I, \S\S6--7]{QuasiRH}; the calculation here is over the fixed field
\(F\) and uses its sextic characters.

Keep the field \(F\), the set \(S\), the target \(\eta\), and the allowed
generators fixed as in \Cref{sec:arithmetic}.  In particular
\(\vartheta=\overline\Lambda\,\overline G\,\eta\), and every character power
is zero on nonunits, including the zeroth power.  The scale notation
\(a,b,M_0,h_0,X,Y\) is fixed in \Cref{eq:reflection-scales}.
All implied constants in this and the next section may depend on the fixed
\(F,S,\eta\) and the fixed weights.  Exponents displayed as numerical
constants do not have that dependence.  An arbitrary positive power loss
may always be chosen smaller before an estimate is applied.

\subsection{Weights and the estimate from reflection}

Choose a nonzero nonnegative function \(W_1\in C_c^\infty((0,\infty))\),
and set
\[
 \widehat W_1(w)=\int_0^\infty W_1(t)t^w\,\frac{dt}{t}.
\]
Choose a tensor Bruhat--Schwartz function \(W_0\) on \(F_S\), compactly
supported in \(F_S^\times\), such that
\begin{equation}
 \widehat W_0(0)=0,\qquad
 M(1/6)\ne0,\qquad
 M(z)=\int_{F_S^\times}\widehat W_0(y)|y|_S^z\,d^\times y .
\label{eq:poisson-Fourier-weight}
\end{equation}
The condition \(\widehat W_0(0)=0\) will remove the zero additive
frequency after Poisson summation.  The condition \(M(1/6)\ne0\) will
ensure that the weight does not kill the residue at \(z=1/6\) in that
expansion.  These are separate from the support condition on \(W_0\),
which removes the original row \(m=0\) from the probe.
Here the Fourier sign is the one fixed in \Cref{sec:arithmetic}.  We
explain the choice and the analytic properties that will be needed.  At
each place the local zeta functional equation gives, for \(0<s<1\),
\[
 \int_{F_v^\times}\widehat W_v(y)|y|_v^s\,d^\times y
   =c_v(s)\int_{F_v^\times}W_v(y)|y|_v^{1-s}\,d^\times y,
 \qquad c_v(1/6)\ne0;
\]
the nonzero constant incorporates the chosen local measures and Fourier
sign \cite[Theorem~2.4.1]{Tate67}.  At one complex place take a smooth annular function
whose additive integral is zero but whose pairing with
\(|y|_v^{5/6}d^\times y\) is nonzero.  Such a function is obtained by
subtracting suitable multiples of two smooth bumps supported in different
radial subannuli: the additive integral is a constant multiple of the
pairing with \(|y|_v d^\times y\), and the two radial weights are not
proportional.  At all other places take nonnegative nonzero annular
functions.  The local identity then proves both requirements in
\Cref{eq:poisson-Fourier-weight}.

For every \(\sigma>0\), each local integral of
\(|\widehat W_v(y)||y|_v^\sigma\) is finite.  At a finite place this follows
by summing the geometric series of shells near zero and using compact
support at infinity; at a complex place it follows from boundedness at
zero and Schwartz decay at infinity.  The same assertions hold after any
number of derivatives in a complex logarithmic norm.  Thus \(M\) is
holomorphic for \(\re z>0\), and integration by parts in that norm gives,
for every compact interval \(J\subset(0,\infty)\) and every \(A>0\),
\begin{equation}
 |M(\sigma+i\tau)|\ll_{J,A}(1+|\tau|)^{-A}
 \quad(\sigma\in J).
\label{eq:poisson-M-decay}
\end{equation}
The analogous estimate holds for \(\widehat W_1\) on every fixed real
strip.

For an exterior ideal \(\mathfrak s\), with an allowed generator \(s\), put
\[
 g_{\mathfrak s}(k)
   =\sum_{d\bmod\mathfrak s}\chi_{\mathfrak s}(d)\e(kd/s).
\]
The value is independent of the allowed generator: a change by
\(\epsilon^6\), \(\epsilon\in U_S\), is absorbed by the same change of
the residue variable.  Use product multiplicative Haar measures for which
the norm decomposition of \(F_S^\times\) has the form
\[
 d^\times y=\frac{dT}{T}\,dt,\qquad y=t\ell_T,\quad t\in F_S^1.
\]
The measure on \(F_S^1/U_S^6\) below is the quotient measure induced by
\(dt\).  This choice also fixes the normalization in \(M\).
Define
\begin{equation}
 \begin{split}
 I(X,Y,Z)=\int_{F_S^1/U_S^6}
 \frac1Y\sum_{\mathfrak s}W_1(q_{\mathfrak s}/Y)
 \frac1{\sqrt{q_{\mathfrak s}X}}\sum_{m\in R}
  q_{\mathfrak s}^{-1/2}g_{\mathfrak s}(-m)
  W_0\!\left(\frac{m}{t_0\ell_{q_{\mathfrak s}X}}\right)
  B_m(Z)\,dt_0 .
 \end{split}
\label{eq:poisson-probe}
\end{equation}
All undecorated ideal sums in this section are over nonzero integral
exterior ideals.  The summand for \(m=0\) vanishes because \(W_0\) is
supported in \(F_S^\times\).  For \(\epsilon\in U_S^6\),
\(\chi_{\mathfrak A}(\epsilon)=\chi_{\mathfrak s}(\epsilon)=1\);
changing \(m\) to \(\epsilon m\) shows that the integrand descends to the
quotient.  That quotient is compact by
\Cref{lem:balanced-generators}, and we use representatives in a fixed
bounded subset of \(F_S^1\).  For each such representative the \(m\)-sum
has finite support.  The Gaussian weight in \(B_m\) makes its ideal sum
absolutely convergent.

We first record the additive estimate in the form also used later for
conductor characters.  If \(\mathfrak s=R\), its one residue class is
understood to give the single fraction \(0\).

\begin{lemma}[Separated additive fractions]\label{lem:additive-sieve}
Let \(K\) be a fixed compact subset of \(F_S\), and let \(N,V\ge2\).
Choose balanced generators \(s\) for the exterior ideals with
\(q_{\mathfrak s}\le V\).  Let \(\mathcal P\) be any subset of the distinct
points
\[
 \left\{\frac d s+R:
       q_{\mathfrak s}\le V,\ d\in(R/\mathfrak s)^\times\right\}
       \subset F_S/R.
\]
For any complex \(b_\alpha\),
\begin{equation}
 \sum_{m\in R\cap\ell_NK}
       \left|\sum_{\alpha\in\mathcal P}b_\alpha\e(m\alpha)\right|^2
 \ll_{F,S,K}(N+V^2)\sum_{\alpha\in\mathcal P}|b_\alpha|^2 .
\label{eq:additive-sieve}
\end{equation}
The dual inequality, with the sums over \(m\) and \(\alpha\) interchanged,
has the same constant.
\end{lemma}

\begin{proof}
Reduced fractions with distinct denominator ideals are distinct modulo
\(R\).  Indeed, at a prime dividing its denominator the valuation of a
reduced fraction is the negative of the denominator valuation; at any
other exterior prime it is nonnegative.  If the denominator valuations
differ, their difference therefore has negative valuation there.
If all agree, congruence modulo \(R\) says the residue classes agree.
For distinct fractions and any \(r\in R\), the numerator in
\[
 \frac d s-\frac {d'}{s'}+r
       =\frac{ds'-d's+rss'}{ss'}
\]
is a nonzero element of \(R\).  The norm identity of
\Cref{lem:s-lattice} gives
\begin{equation}
 \left|\frac d s-\frac {d'}{s'}+r\right|_S
       \ge (q_{\mathfrak s}q_{\mathfrak s'})^{-1}\ge V^{-2}.
\label{eq:additive-separation}
\end{equation}
For \(r\in R\setminus\{0\}\) it also gives \(|r|_S\ge1\).

Put \(d_F=[F:\Q]\) and
\(\delta=c\min(N^{-1/d_F},V^{-2/d_F})\), with \(c>0\) fixed sufficiently
small.  There is a nonnegative Bruhat--Schwartz function \(\Phi\) which is
at least one on \(\ell_NK\), has
\(\widehat\Phi(0)\ll\delta^{-d_F}\), and whose Fourier transform is supported
in a product of fixed finite-place additive subgroups and complex disks
of ordinary radius \(C\delta\).  Here is a direct construction.  At each
complex place take the inverse transform of a smooth bump
supported in a disk, evaluated at \(\delta y_v\).  Its value stays
bounded below when \(\delta|y_v|\) is sufficiently small.  At finite
places use the inverse transform of the indicator of a sufficiently
small compact open subgroup, whose annihilator contains the finite
projection of \(K\).  The product, multiplied by a fixed constant and
squared in modulus, majorizes the box.  Fourier support follows from
convolution, and a change of scale gives integral \(O(\delta^{-d_F})\).
Shrinking \(c\) makes this support contain no translate by \(R\) of a
difference of distinct points in \(\mathcal P\), by
\Cref{eq:additive-separation}: the product of the complex squared absolute
values on that support is \(O(\delta^{d_F})\), and its finite-place product
is bounded by a fixed constant.  It contains no nonzero point of \(R\)
for the same reason.

Expand the square after majorizing it by \(\Phi\), and apply Poisson
summation on \(R\).  For \(\alpha,\beta\in\mathcal P\),
\[
 \sum_{m\in R}\Phi(m)\e(m(\alpha-\beta))
   =\sum_{r\in R}\widehat\Phi(r-\alpha+\beta)
   =\one_{\alpha=\beta}\widehat\Phi(0).
\]
Since \(\delta^{-d_F}\ll N+V^2\), this proves
\Cref{eq:additive-sieve}.  Equality of the operator norms of a finite
matrix and its adjoint gives the dual assertion.
\end{proof}

\begin{proposition}\label{prop:probe-low}
For the weights above, for every \(\epsilon>0\) and \(Z\ge2\),
\begin{equation}
 I(Z^a,Z^b,Z)\ll_{F,S,\eta,W_0,W_1,\epsilon}
 Z^{(2M_0-1-b)/2+\epsilon}
       =Z^{933/2000+\epsilon}.
\label{eq:probe-low}
\end{equation}
\end{proposition}

\begin{proof}
Write \(v=q_{\mathfrak s}/Y\) and
\(y=m/(t_0\ell_{XY})\).  For \(v\) in the support of \(W_1\), the support
of \(W_0(y/\ell_v)\) puts \(y\) in one fixed compact annular set.
On this set the smooth function
\[
 v^{-1/2}W_1(v)W_0(y/\ell_v)
\]
has a Fourier series in \(\log v\) whose coefficients \(c_j(y)\)
satisfy \(\sum_j\|c_j\|_\infty<\infty\), uniformly in \(y\).
To see this, take a compact interval containing the logarithmic support,
extend by zero to a slightly larger period, and integrate twice by parts.
All the required derivatives are bounded uniformly on the compact
\(y\)-set.  Thus we can separate this weight from the additive sums with
a bounded total coefficient cost.

For each Fourier term, expansion of \(g_{\mathfrak s}\) gives a
coefficient of modulus \(Y^{-1}q_{\mathfrak s}^{-1/2}\) at each reduced
fraction \(d/s\), with \(q_{\mathfrak s}\asymp Y\).  Their total squared
mass is at most
\[
 Y^{-2}\sum_{q_{\mathfrak s}\asymp Y}
                 q_{\mathfrak s}^{-1}|(R/\mathfrak s)^\times|
 \ll_F Y^{-1},
\]
using the fixed-field ideal count.  Take \(N=XY\) and \(V\ll Y\) in
\Cref{lem:additive-sieve}.  The bounded set of representatives \(t_0\)
only enlarges the fixed compact row box.  By the triangle inequality for
the Fourier series, the squared \(\ell^2\) norm in \(m\) of the whole
additive factor in \Cref{eq:poisson-probe}, including
\((q_{\mathfrak s}X)^{-1/2}\), is
\[
 \ll (XY)^{-1}\frac{XY+Y^2}{Y}\ll Y^{-1}.
\]
The last inequality uses \(b<a\).  On the same enlarged compact annular
row box, \Cref{prop:reflection} gives
\[
 \sum_m |B_m(Z)|^2\ll_\epsilon Z^{2M_0-1+\epsilon}.
\]
Cauchy--Schwarz in \(m\), followed by integration over the compact
quotient, proves \Cref{eq:probe-low}, after decreasing the preliminary
loss.
\end{proof}

\subsection{Poisson summation and three Mellin variables}

For an exterior ideal \(\mathfrak A\) with allowed generator \(A\), write
\[
 g_{\chi_{\mathfrak A}}(A,k)
   =\sum_{v\bmod\mathfrak A}\chi_{\mathfrak A}(v)\e(kv/A)
\]
and, for \(h\in R\), define
\begin{equation}
 F(s,A,h)=
 \sum_{\substack{d\bmod\mathfrak s\\h\equiv Ad\pmod{\mathfrak s}}}
 \chi_{\mathfrak s}(d)\,
 g_{\chi_{\mathfrak A}}\!\left(A,\frac{h-Ad}{s}\right).
\label{eq:finite-Poisson-factor}
\end{equation}
The argument of the last Gauss sum is in \(R\) by the congruence.  The
definitions for the unit ideal mean a single residue and the value \(1\)
for its character.

The finite Fourier transform of the row coefficient modulo
\(\mathfrak s\mathfrak A=sAR\) is exactly
\begin{equation}
 \sum_{m\bmod\mathfrak s\mathfrak A}
 q_{\mathfrak s}^{-1/2}g_{\mathfrak s}(-m)
 \chi_{\mathfrak A}(m)\e(hm/(sA))
       =q_{\mathfrak s}^{1/2}F(s,A,h).
\label{eq:finite-Poisson-transform}
\end{equation}
In fact expand \(g_{\mathfrak s}\), write \(m=v+Aj\) with
\(v\bmod\mathfrak A\) and \(j\bmod\mathfrak s\), and sum over \(j\).
The sum of \(\e((h-Ad)j/s)\) is \(q_{\mathfrak s}\) precisely when
\(h\equiv Ad\pmod{\mathfrak s}\), and is zero otherwise.  The remaining
sum over \(v\) is the one in \Cref{eq:finite-Poisson-factor}.  This
argument does not require \(\mathfrak s\) and \(\mathfrak A\) to be
coprime.

The dilation \(t_0\ell_{q_{\mathfrak s}X}\) has additive modulus
\(q_{\mathfrak s}X\), and the lattice \(sAR\) has index
\(q_{\mathfrak s}q_{\mathfrak A}\).  Apply Poisson summation to the
\(m\)-sum and use \Cref{eq:finite-Poisson-transform}.  Its dilation divided
by the index is \(X/q_{\mathfrak A}\); after the two square root factors
in the probe the multiplier is
\[
 (q_{\mathfrak s}X)^{-1/2}
       \frac X{q_{\mathfrak A}}q_{\mathfrak s}^{1/2}
       =\frac{X^{1/2}}{q_{\mathfrak A}}.
\]
The Fourier argument at \(h\) is
\[
 t_0\ell_{q_{\mathfrak s}X}\frac h{sA},\qquad
 \left|\ell_{q_{\mathfrak s}X}\frac h{sA}\right|_S
       =\frac{Xq_h}{q_{\mathfrak A}}.
\]
These norm factors will determine all three Mellin exponents.

For \(\epsilon\in U_S^6\), changing \(d\) to \(\epsilon d\) and the
Gauss residue variable \(v\) to \(\epsilon^{-1}v\) in
\Cref{eq:finite-Poisson-factor} shows
\(F(s,A,\epsilon h)=F(s,A,h)\).  Sum the nonzero \(h\)'s in such
orbits and unfold the quotient \(F_S^1/U_S^6\).  If
\[
 V_0(T)=\int_{F_S^1}\widehat W_0(t\ell_T)\,dt,
\]
the orbit of \(h\) contributes
\(F(s,A,h)V_0(Xq_h/q_{\mathfrak A})\).  Multiplication by the norm-one
part of \(\ell_{q_{\mathfrak s}X}h/(sA)\) preserves the sheet measure, so
no shape factor remains.  This unfolding is absolutely justified: for
each \(\sigma>0\), the local integrability used above gives an absolute
sheet majorant \(O_\sigma(T^{-\sigma})\), by using one complex weighted
supremum and the other weighted local integrals.  Taking \(\sigma>1\)
sums the \(q_u^{-\sigma}q_{\mathfrak r}^{-6\sigma}\) majorant, while
the Gaussian handles the ideal \(\mathfrak A\); this permits Tonelli's
theorem in the orbit sum.  The zero frequency contributes
nothing because \(\widehat W_0(0)=0\).

By the measure disintegration, the Mellin transform of \(V_0\) is \(M\).
The same weighted derivative majorants in the complex norm variables
justify differentiation of the sheet integral in \(\log T\), so
\(V_0(T)\) is smooth for \(T>0\).  Its absolute weighted integrability
and \Cref{eq:poisson-M-decay} then give pointwise Mellin inversion by
Fourier inversion in \(\log T\) on each line \(\re z>0\).
By \Cref{eq:frequency-decomposition}, write the orbits uniquely as
\(h=ur^6\), where \(u\in\mathcal U\) and \(\mathfrak r\) is an exterior
ideal with allowed generator \(r\).  The principal row \(u=1\) consists
of the surviving nonzero sixth-power frequencies \(h=r^6\).  Hence the
precise Poisson expression is
\begin{equation}
 \begin{split}
 I(X,Y,Z)
  ={}&\frac{X^{1/2}}Y
   \sum_{\substack{\mathfrak c\ {\rm squarefree}\\
                    \mathfrak n,\mathfrak s}}
   \frac{\gamma_2(\mathfrak c)\vartheta(\mathfrak A)}
        {q_{\mathfrak c}^{1/2}q_{\mathfrak n}q_{\mathfrak A}}\,
   P_G(q_{\mathfrak A}/Z)W_1(q_{\mathfrak s}/Y)\\
  &{}\times
   \sum_{u\in\mathcal U}\sum_{\mathfrak r}
     F(s,A,ur^6)
     V_0\!\left(\frac{Xq_uq_{\mathfrak r}^6}{q_{\mathfrak A}}\right),
 \qquad \mathfrak A=\mathfrak c\mathfrak n^3.
 \end{split}
\label{eq:Poisson-orbits}
\end{equation}
Here \(P_G\) is the Gaussian of \Cref{eq:reflection-gaussian}.

Use Mellin inversion for \(P_G,W_1,V_0\), with transforms
\(e^{t^2},\widehat W_1(w),M(z)\), respectively.  The factors involving
\(\mathfrak A\) are
\(q_{\mathfrak A}^{-1}q_{\mathfrak A}^{-t}q_{\mathfrak A}^{z}\).
Set \(t=\xi+z-1\).  We obtain
\begin{equation}
 \begin{split}
 I(X,Y,Z)=\frac1{(2\pi i)^3}
  \int_{(4)}\int_{(4)}\int_{(2)}
  &X^{1/2-z}Z^{\xi+z-1}Y^{w-1}
   e^{(\xi+z-1)^2}M(z)\widehat W_1(w)\\
  &{}\times\sum_{u\in\mathcal U}q_u^{-z}
       \mathcal F_u(\xi,w,z)\,dz\,dw\,d\xi ,
 \end{split}
\label{eq:poisson-mellin}
\end{equation}
where the innermost line is \(\re z=2\), the middle line is
\(\re w=4\), and the outer line is \(\re\xi=4\), and
\begin{equation}
 \mathcal F_u(\xi,w,z)=
 \sum_{\substack{\mathfrak c\ {\rm squarefree}\\
                  \mathfrak n,\mathfrak s,\mathfrak r}}
 \gamma_2(\mathfrak c)\vartheta(\mathfrak A)F(s,A,ur^6)
 q_{\mathfrak c}^{-1/2-\xi}q_{\mathfrak n}^{-1-3\xi}
 q_{\mathfrak s}^{-w}q_{\mathfrak r}^{-6z},
 \qquad \mathfrak A=\mathfrak c\mathfrak n^3.
\label{eq:calF-definition}
\end{equation}
All interchanges used for \Cref{eq:poisson-mellin} are in absolute
convergence.  For example,
\(|F(s,A,h)|\le q_{\mathfrak s}q_{\mathfrak A}\),
\(|\gamma_2(\mathfrak c)|=1\), and the ideal count, including the finitely
many S-unit classes for \(u\), give absolute convergence on these lines.
The factors \(M(z)\) and \(\widehat W_1(w)\) have arbitrary vertical
polynomial decay, while the Gaussian has Gaussian decay in
\(\im(\xi+z)\).  Changing height variables to
\((\im(\xi+z),\im z,\im w)\) proves absolute integrability as well.

\subsection{The local correction}

We extend the notation \(L^S(s,\eta)\) from
\Cref{sec:coefficient-family} to a Hecke character on exterior primes:
the factors at \(S\) are omitted, and any additional zero prime values
are retained in the Euler product.
Thus, for example,
\[
 L^S(w,\chi_\bullet(u))
    =\prod_{\mathfrak p\notin S}
       (1-\chi_{\mathfrak p}(u)q_{\mathfrak p}^{-w})^{-1}
 \quad(\re w>1).
\]
As before,
\(\zeta_F^S(s)=\prod_{\mathfrak p\notin S}(1-q_{\mathfrak p}^{-s})^{-1}\).

\begin{proposition}\label{prop:euler-product}
For every \(u\in\mathcal U\), the series in
\Cref{eq:calF-definition} has the Euler factorization
\begin{equation}
 \mathcal F_u(\xi,w,z)=
 \frac{\zeta_F^S(6z)L^S(w,\chi_\bullet(u))}
      {L^S(\xi,\eta\overline{\chi_\bullet(u)})}
             \mathcal H_u(\xi,w,z)
\label{eq:probe-euler-product}
\end{equation}
initially in absolute convergence.  The correction \(\mathcal H_u\) is
holomorphic in open neighborhoods of both real-part regions
\begin{equation}
 \begin{array}{lll}
 \re\xi\ge .998,&\re w\ge .003,&\re z\ge .4,\\
 \re\xi\ge .998,&\re w\ge .95,&\re z\ge .16.
 \end{array}
\label{eq:euler-regions}
\end{equation}
For every \(\epsilon>0\) it satisfies
\[
 |\mathcal H_u(\xi,w,z)|\ll_{F,S,\eta,\epsilon}q_u^\epsilon
\]
there, uniformly in all heights.  If \(S\) contains every prime of norm
at most a sufficiently large fixed constant, then
\(\mathcal H_1\) is bounded and nowhere zero on both neighborhoods.
\end{proposition}

The two regions serve different contour movements in the next section.
The first permits the small and intermediate nonprincipal rows to place
\(w\) at \(.003\) and \(z\) at \(.4\), while \(\xi\) can reach \(.998\)
once the reciprocal denominator has been controlled.  The second contains
the principal numerator poles \(w=1\), \(z=1/6\) and the nearby lines
\(w=.95\), \(z=.16\) used to isolate their double residue.  The
nowhere-zero conclusion for \(u=1\) will ensure that its correction does
not cancel a singularity of \(1/L^S(\xi,\eta)\).

\begin{proof}
We first factor the finite expression, keeping the units that Chinese
remaindering introduces.  Use products of allowed prime generators for
\(s,A,r\).  They differ from any other allowed choices by sixth powers
of S-units, which have no effect on the expressions.  At an exterior
prime \(\mathfrak p\), with generator \(p\) and norm \(Q=q_{\mathfrak p}\),
write
\[
 \begin{gathered}
 t=v_{\mathfrak p}(\mathfrak A)=e_0+3l,\quad
 e_0=v_{\mathfrak p}(\mathfrak c)\in\{0,1\},\quad
 l=v_{\mathfrak p}(\mathfrak n),\quad k=v_{\mathfrak p}(\mathfrak s),\\
 j=v_{\mathfrak p}(u),\quad m'=v_{\mathfrak p}(\mathfrak r),
 \quad J=j+6m'.
 \end{gathered}
\]
Write \(s=p^ks_p\), \(A=p^tA_p\), and \(ur^6=p^Jh_p\); the subscripts
mean that the full \(p\)-power has been omitted.  These three subscripted
quantities are units at \(\mathfrak p\).  Put
\(\psi_{\mathfrak p}^{\circ}(x)=\psi_{\mathfrak p}(-x)\).
For \(x\in F\), the global additive product gives
\(\e(x)=\prod_{\mathfrak p\notin S}\psi_{\mathfrak p}^{\circ}(x)\),
with only finitely many nontrivial factors.  Expand the inner Gauss sum
in \Cref{eq:finite-Poisson-factor} and apply the Chinese remainder
decomposition to each residue ring.  Before changing variables, the local
factor of \(F(s,A,ur^6)\) is
\[
 \sum_{\substack{d\bmod p^k\\p^Jh_p\equiv p^tA_p d\pmod{p^k}}}
 \chi_{\mathfrak p}(d)^k
 \sum_{v\bmod p^t}\chi_{\mathfrak p}(v)^t
 \psi_{\mathfrak p}^{\circ}\!\left(
  \frac{(p^Jh_p-p^tA_p d)v}{p^{k+t}s_pA_p}\right).
\]
Here an absent modulus \(p^0\) has one residue and character value \(1\).
For a positive modulus the character retains its unit mask, even when
its exponent is divisible by six.  Now make the changes
\[
 d=(h_p/A_p)d',\qquad v=(s_pA_p/h_p)v'.
\]
They are bijections of the relevant residue rings.  The congruence becomes
\(p^J\equiv p^t d'\pmod{p^k}\), and the local additive argument becomes
\((p^J-p^td')v'/p^{k+t}\).  The unit multiplier removed by these changes
is exactly
\begin{equation}
 \chi_{\mathfrak p}(h_p)^{k-t}
 \chi_{\mathfrak p}(s_p)^t
 \chi_{\mathfrak p}(A_p)^{t-k}.
\label{eq:Euler-CRT-unit}
\end{equation}
There are no hidden coprimality assumptions here: these changes only use
the units left after omitting the full \(p\)-powers.  Moreover
\(\chi_{\mathfrak p}(h_p)=\chi_{\mathfrak p}(u/p^j)\), because the
remaining factors of \(r^6\) are sixth powers.

For completeness define the stripped local sum by
\begin{equation}
 C_{t,k}(J)=
 \sum_{\substack{d\bmod p^k\\p^J\equiv p^td\pmod{p^k}}}
 \chi_{\mathfrak p}(d)^k
 \sum_{v\bmod p^t}\chi_{\mathfrak p}(v)^t
       \psi_{\mathfrak p}^{\circ}
             \left(\frac{(p^J-p^td)v}{p^{k+t}}\right).
\label{eq:C-local-definition}
\end{equation}
When \(k=0\) the outer modulus is one, its character is one, and its
single residue may be taken to be zero.  The analogous convention applies
to \(t=0\) in the inner sum.  When \(k>0\) or \(t>0\), the corresponding
character is supported on units even if its exponent is divisible by
six.  These conventions distinguish the absent modulus from the zeroth
power character on a nontrivial modulus.

We check explicitly that the remaining factors in different primes
separate.  Expand \Cref{eq:Euler-CRT-unit} using the product generators,
and consider an unordered pair of distinct primes \(\mathfrak p,\mathfrak
p'\).  In their common symbol
\(\chi_{\mathfrak p}(p')=\chi_{\mathfrak p'}(p)\), the contributions from
\(\chi_{\mathfrak p}(s_p)^t\) and
\(\chi_{\mathfrak p}(A_p)^{t-k}\), together with the analogous factors at
\(\mathfrak p'\), have exponent
\[
 tk'+t'k+(t-k)t'+(t'-k')t=2tt'.
\]
In particular all \(s\)-\(A\) interactions cancel.  By
\Cref{eq:gauss-crt}, the Chinese remainder factors in
\(\gamma_2(\mathfrak c)\) add \(4e_0e_0'\).  Their total is
\[
 2(e_0+3l)(e_0'+3l')+4e_0e_0'\equiv0\pmod6.
\]
Thus all cross-prime terms cancel, including when both moduli have powers
of the same primes; those shared powers were retained inside
\Cref{eq:C-local-definition}.  The \(h_p\)-factors remain precisely as
characters of \(u/p^j\).  If
\(\theta_{\mathfrak p}=\chi_{\mathfrak p}(u/p^j)\in\mu_6\), the exact
local factor is consequently
\begin{equation}
 \begin{split}
 \mathcal F_{u,\mathfrak p}(\xi,w,z)
  =\sum_{\substack{e_0=0,1\\l,k,m'\ge0}}
  &\gamma_2(\mathfrak p)^{e_0}\vartheta(\mathfrak p)^{e_0+3l}
   \theta_{\mathfrak p}^{\,k-e_0-3l}C_{e_0+3l,k}(j+6m')\\
  &{}\times
  Q^{-e_0(1/2+\xi)-l(1+3\xi)-kw-6m'z}.
 \end{split}
\label{eq:exact-Euler-factor}
\end{equation}
It follows initially by absolute convergence that
\(\mathcal F_u=\prod_{\mathfrak p\notin S}\mathcal F_{u,\mathfrak p}\).

Here are explicit evaluations of every \(C_{t,k}(J)\).  Put
\[
 T_v(y)=\sum_{x\bmod p}\chi_{\mathfrak p}(x)^v
                   \psi_{\mathfrak p}^{\circ}(yx/p),
 \qquad \tau_v=T_v(1).
\]
The unit mask is retained for \(v\equiv0\pmod6\).  Thus
\[
 T_v(y)=
 \begin{cases}
  \overline{\chi_{\mathfrak p}(y)}^{\,v}\tau_v,&p\nmid y,\\
  Q-1,&p\mid y,\ v\equiv0\pmod6,\\
  0,&p\mid y,\ v\not\equiv0\pmod6,
 \end{cases}
 \qquad
 |\tau_v|=\sqrt Q\ (v\not\equiv0\!\!\pmod6),\quad \tau_0=-1.
\]
For \(t=0\) the value is \(C_{0,0}(J)=1\), and
\(C_{0,k}(J)=\one_{J=0}\) for \(k>0\).  For \(t>0\) it is
\begin{equation}
 \begin{split}
 C_{t,0}(J)&=\one_{J\ge t-1}Q^{t-1}T_t(p^{J-t+1}),\\
 C_{t,1}(J)&=\one_{J\ge t}\tau_1Q^{t-1}T_{t-1}(p^{J-t}),\\
 C_{t,k}(J)&=0\qquad(k\ge2).
 \end{split}
\label{eq:C-local-values}
\end{equation}
For the first formula sum the \(Q^{t-1}\) lifts of a residue \(v\bmod p\)
in the inner sum; they cancel unless \(J\ge t-1\).  For \(k=1\) the
congruence first requires \(J\ge1\).  Sum \(d\bmod p\) before \(v\).
Because \(v\) is a unit, the sum is
\(\tau_1\overline{\chi_{\mathfrak p}(v)}\);
\(\chi_{\mathfrak p}(-1)=1\) removes the sign.  The remaining sum over
lifts of \(v\) is zero unless \(J\ge t\), and then gives the second
formula.  For \(k\ge2\), varying \(d\) by \(p^{k-1}\lambda\), \(\lambda\bmod p\),
preserves the congruence and its unit character.  For every unit \(v\)
the phase varies by \(\psi_{\mathfrak p}^{\circ}(-\lambda v/p)\), whose sum in
\(\lambda\) is zero.  This proves the last formula.  In particular, for \(t>0\),
\begin{equation}
 |C_{t,0}(J)|\le Q^t,\qquad
 |C_{t,1}(J)|\le Q^{t+1/2},
\label{eq:C-local-bounds}
\end{equation}
with the support restrictions visible in \Cref{eq:C-local-values}.

Suppose first that \(\mathfrak p\nmid(u)R\), so \(j=0\).  Set
\[
 V=Q^{-6z},\quad
 W=\chi_{\mathfrak p}(u)Q^{-w},\quad
 D_1=\eta(\mathfrak p)\overline{\chi_{\mathfrak p}(u)}Q^{-\xi},
 \quad
 E_1=\vartheta(\mathfrak p)^6Q^{4-6\xi}V.
\]
The \(t=0\) values give
\((1-V)^{-1}+W/(1-W)\): a positive \(k\) requires \(J=0\).
The first squarefree theta term has
\(e_0=1\), \(l=k=m'=0\), hence \(t=1\), \(J=0\), and
\(C_{1,0}(0)=\tau_1\).  Its contribution is
\[
 \begin{split}
 &\gamma_2(\mathfrak p)\vartheta(\mathfrak p)
  \overline{\chi_{\mathfrak p}(u)}Q^{-1/2-\xi}\tau_1\\
 &\qquad=\gamma_1(\mathfrak p)\gamma_2(\mathfrak p)
   \vartheta(\mathfrak p)\overline{\chi_{\mathfrak p}(u)}Q^{-\xi}
   =-D_1.
 \end{split}
\]
Indeed \(\tau_1=\sqrt Q\,\gamma_1(\mathfrak p)\) in the chosen additive
sign, and \Cref{eq:gauss-corrections} gives
\(\gamma_1\gamma_2\vartheta=-\eta\).  This negative term is the first
nonconstant term of the reciprocal local \(L\)-factor \(1-D_1\); it is
the reason that \(L^S(\xi,\eta\overline{\chi_\bullet(u)})\) occurs in the
denominator of \Cref{eq:probe-euler-product}.  We extract
\((1-V)^{-1}(1-W)^{-1}(1-D_1)\), whose first two factors are those of
\(\zeta_F^S(6z)\) and \(L^S(w,\chi_\bullet(u))\).  The mixed terms and
the higher theta terms will be included in the correction.

For \(t>0\) and \(J=6m'\), the formulas show that only
\(t\equiv0,1\pmod6\) occur.  For \(t=6d\), \(d\ge1\), the \(k=0\) sum is
\[
 E_1^d\frac{1-Q^{-1}}{1-V},
\]
and the \(k=1\) term is \(E_1^dW\).  For the latter, the only possible
valuation is \(J=t\), and \(\tau_1\tau_5=Q\).  For \(t=6d+1\), \(d\ge0\),
the \(k=0\) term is \(-D_1E_1^d\), extending the first squarefree term
just computed.  Finally its \(k=1\) sum is
\[
 -(Q-1)D_1E_1^d\,\frac{WV}{1-V}.
\]
Here \(J>t\) and \(T_0(p^{J-t})=Q-1\); the first possible extra
\(m'\) gives the displayed \(V\).  Summing in \(d\) proves the closed
formula
\begin{equation}
 \begin{split}
 \mathcal F_{u,\mathfrak p}
  ={}&\frac1{1-V}+\frac W{1-W}
   +\frac{E_1}{1-E_1}
           \left(\frac{1-Q^{-1}}{1-V}+W\right)\\
   &-\frac{D_1}{1-E_1}
           \left(1+\frac{(Q-1)WV}{1-V}\right).
 \end{split}
\label{eq:Euler-good-exact}
\end{equation}
This derivation also explains the vanishing of the allowed theta exponents
\(t\equiv3,4\pmod6\): the possible boundary values of \(J\) in
\Cref{eq:C-local-values} have the wrong residue modulo six, and the
remaining \(T\)-sum is a nontrivial character sum at zero.

Define at every prime
\[
 \mathcal H_{u,\mathfrak p}
   =(1-V)\frac{1-W}{1-D_1}\mathcal F_{u,\mathfrak p},
\]
where \(W=D_1=0\) at \(\mathfrak p\mid(u)R\), in accordance with the
zero extension.  We verify normal convergence on explicit open
neighborhoods: in each line of \Cref{eq:euler-regions} decrease all three
lower endpoints by \(10^{-5}\), and use the resulting strict inequalities.
On these neighborhoods all of \(V,W,D_1,E_1\) have modulus bounded away
from one once \(Q\) exceeds a fixed constant.  Subtracting \(1\) after
multiplying \Cref{eq:Euler-good-exact} by
\((1-V)(1-W)/(1-D_1)\) gives
\begin{equation}
 |\mathcal H_{u,\mathfrak p}-1|
 \ll |VW|+|D_1V|+|D_1W|+Q|D_1WV|+|E_1|
 \ll Q^{-1-\delta},\qquad \delta=\frac1{4000},
\label{eq:Euler-good-error}
\end{equation}
uniformly for \(\mathfrak p\nmid(u)R\).  For clarity, before the
\(10^{-5}\) enlargement the weakest of the first three terms is
\(|D_1W|\le Q^{-1.001}\), in the first region.  The exponents in the
last two terms are at most
\[
 1-\re\xi-\re w-6\re z\le-1.908,\qquad
 4-6\re\xi-6\re z\le-2.948
\]
on the union of the two regions.  The enlargement leaves respectively
the margins \(0.00098\), \(0.90792\), and \(1.94788\) beyond exponent
\(-1\), which suffice for the stated \(\delta\).  More explicitly,
the specialization at \(E_1=0\) satisfies
\[
 \left.\mathcal H_{u,\mathfrak p}\right|_{E_1=0}-1
 =\frac{D_1(W+V-WV)-WV-(Q-1)D_1(1-W)WV}{1-D_1}.
\]
This accounts for the \(Q|D_1WV|\) term as well as the first three
terms in \Cref{eq:Euler-good-error}.  The difference from this
specialization is \(O(|E_1|)\) on the stated neighborhoods: the remaining
denominators are bounded, and \(Q|D_1WV|<1\) once \(Q\) exceeds the fixed
cutoff.

At \(\mathfrak p\mid(u)R\), \(1\le j\le5\).  The \(t=0\) contribution
is \((1-V)^{-1}\).  The term \(t=1,k=0\) vanishes, since \(J\ge1\) and
\(T_1(p^J)=0\).  By \Cref{eq:C-local-bounds}, the sum for \(t=1,k=1\)
has modulus at most
\[
 \frac{Q^{1-\re\xi-\re w}}{1-|V|}.
\]
Its exponent is at most \(-0.001\) on the first region and at most
\(-0.948\) on the second.  For \(l\ge1\), before the factor \(V^{m'}\),
the bound for the summand with \(e_0\in\{0,1\}\) is at most a constant
times \(Q\) to the power
\begin{equation}
 (2-3\re\xi)l+(1/2-\re\xi)e_0
                    +\max(0,1/2-\re w).
\label{eq:Euler-bad-exponents}
\end{equation}
This is at most \(-0.497\) for \(l=1,e_0=0\) on the first region and
decreases by at least \(0.994\) for each further \(l\); the second region
has a larger saving.  The \(10^{-5}\) enlargement still leaves a saving
greater than \(0.49\) in this sum, and leaves \(0.00098\) in the
\(t=1,k=1\) term.  The geometric sums in \(l,m'\), followed by
multiplication by \(1-V\), therefore give
\begin{equation}
 \mathcal H_{u,\mathfrak p}=1+O(Q^{-\delta})
 \quad(\mathfrak p\mid(u)R),\qquad \delta=\frac1{4000},
\label{eq:Euler-bad-error}
\end{equation}
on the same neighborhoods.  These estimates prove absolute local
convergence there as well.

The product of the factors in \Cref{eq:Euler-good-error} converges
normally, since \(\sum_{\mathfrak p}q_{\mathfrak p}^{-1-\delta}<\infty\).
The remaining primes are precisely the divisors of \((u)R\).  For every
\(\epsilon>0\),
\[
 \prod_{\mathfrak p\mid(u)R}(1+Cq_{\mathfrak p}^{-\delta})
       \ll_{F,S,\epsilon}q_u^\epsilon:
\]
for sufficiently large \(Q\) each local factor is at most \(Q^\epsilon\),
and the finitely many smaller primes contribute a fixed constant.
Consequently \(\mathcal H_u=\prod_{\mathfrak p\notin S}
\mathcal H_{u,\mathfrak p}\) is holomorphic and obeys the asserted bound.
Multiplication of its defining local factors gives
\Cref{eq:probe-euler-product} in the initial region and then as a
meromorphic identity on either neighborhood.

For \(u=1\) every prime satisfies \Cref{eq:Euler-good-error}.  The local
estimate determines a fixed norm cutoff beyond which its uniform bound
is smaller than \(1/2\).  Choose \(S\) to contain all primes up to that
cutoff from the outset, as permitted in \Cref{sec:arithmetic}.  Every remaining
\(\mathcal H_{1,\mathfrak p}\) is nonzero, and its logarithm defined near
\(1\) has an absolutely and uniformly convergent sum.  Its exponential is
\(\mathcal H_1\), which is therefore bounded and nowhere zero on both
neighborhoods.  This proves the proposition.
\end{proof}

For the remainder of the proof, write
\[
 \mathcal H_\eta(s):=\mathcal H_1(s,1,1/6).
\]
This is the principal specialization announced in
\Cref{sec:coefficient-family}.  Since \((w,z)=(1,1/6)\) lies in the
second region of \Cref{eq:euler-regions}, \(\mathcal H_\eta\) is
holomorphic, bounded, and nowhere zero on an open neighborhood of the
closed half-plane \(\re s\ge .998\), uniformly in height.
